\clearpage
\section{Конструкторская часть}

\subsection{Алгоритм полного перебора}

Условные обозначения схем приняты по ГОСТ~19.701-90~\cite{flowstandard}.
Стрелки входят в символы сверху. Парные соединители с одинаковой буквой обозначают продолжение линии в другом месте схемы.

На рисунке~\ref{fig:exact} перестановки портов перебираются в лексикографическом порядке. Для каждой проверяется до $n-1$ переходов. Допустимый путь сравнивается с лучшим по паре $(L,C)$.

\flowchart{exact}{Схема алгоритма полного перебора}{fig:exact}

\clearpage
\subsection{Алгоритм с элитными муравьями}

На рисунке~\ref{fig:ants} показаны $D$ итераций по $K$ попыток. Исходно $\tau_{ij}=1$. В начале итерации матрица добавок $\Delta$ обнуляется.

Допустимый путь с положительным временем $L$ добавляет $Q/L$ на каждое своё ребро. После всех $K$ попыток рёбра лучшего пути с начала поиска получают ещё $eQ/L(P^*)$. Для одноэлементного пути добавок нет. Используются $Q=100$, $e=4$. По окончании итерации феромон обновляется симметрично
\begin{equation}
\tau_{ij}\gets \min\!\bigl(10,\max(0{,}05,(1-\rho)\tau_{ij}+\Delta_{ij})\bigr),
\qquad \tau_{ji}\gets\tau_{ij},
\label{eq:evaporation}
\end{equation}
\hspace*{\parindent}где $\rho\in[0,1]$ --- доля испарения. Муравьи одной итерации используют одну матрицу $\tau$.

\flowchart{ants}{Общая схема муравьиного алгоритма}{fig:ants}

\clearpage
\subsection{Выбор перехода и соблюдение бюджета}

Построение пути показано на рисунке~\ref{fig:walk}. Вес ребра из текущего порта $i$ в непосещённый соседний порт $j\in S_i$ равен
\begin{equation}
w_{ij} = \tau_{ij}^{\alpha}(1/t_{ij})^{\beta}\frac{F}{F+c_{ij}},
\label{eq:weight}
\end{equation}
\hspace*{\parindent}где $\alpha$ и $\beta$ задают влияние феромона и времени, а $F=1+(B-s)/r$ учитывает остаток бюджета на $r$ переходов. Вероятность выбора равна
\begin{equation}
p_{ij}=(1-\varepsilon)\frac{w_{ij}}{\sum_{k\in S_i} w_{ik}}+
       \frac{\varepsilon}{|S_i|},
\qquad \varepsilon=0{,}05.
\label{eq:probability}
\end{equation}

Поэтому для каждого $j\in S_i$ справедлива оценка $p_{ij}\geq\varepsilon/|S_i|>0$.

\flowchart{walk}{Схема построения маршрута одним муравьём}{fig:walk}

\clearpage
\subsection{Допустимость результата и завершение}

Каждый открытый гамильтонов путь представлен одной перестановкой. После проверки очередной перестановки сохраняется лучший допустимый путь среди всех проверенных. При превышении бюджета проверка текущей перестановки прекращается, поскольку все стоимости положительны. Так перебор находит оптимум или устанавливает отсутствие решения.

Кандидаты на переход включают все рёбра к непосещённым портам. Если выбранное ребро превышает остаток бюджета либо соседей нет, муравей завершает попытку без результата и добавки феромона.

На каждом успешном шаге в маршрут добавляется непосещённый сосед, а стоимость остаётся в пределах $B$. За не более чем $n-1$ таких шагов строится путь из $n$ различных портов, удовлетворяющий~\eqref{eq:objective}. За $DK$ попыток эвристика может пропустить все допустимые пути.

Любой допустимый путь может быть построен с положительной вероятностью. Его начальный порт выбирается с вероятностью $1/n$, а вероятность каждого следующего перехода положительна по формуле~\eqref{eq:probability}. По окончании поиска выбирается лучший из построенных путей.

\subsection{Трудоёмкость по схемам алгоритмов}

Используется модель единичной стоимости для арифметики, сравнения, обращения к элементу массива и генерации отдельного случайного числа. Для фиксированных машинных типов вычисление веса перехода также считается операцией постоянной трудоёмкости. Оценки относятся к поиску с уже подготовленной матрицей графа.

По схеме полного перебора (рисунок~\ref{fig:exact}) инициализация требует $\Theta(n)$ операций. Блок оценки маршрута для перестановки $P$ проверяет $\ell(P)$ рёбер, $0\le\ell(P)\le n-1$. Обновление результата копирует $n$ вершин только при улучшении. Генерация следующей перестановки занимает $O(n)$ операций на поиск элемента, обмен и разворот суффикса. Для $U$ фактических улучшений получаем верхнюю оценку
\begin{equation}
F_{\text{п}}(n)\ \leq\ a n + \sum_{P}\bigl(b+c\ell(P)+dn\bigr)+fUn
 = O(n\cdot n!),\quad U\le n!.
\end{equation}
Константы $a,b,c,d,f$ не зависят от $n$. На полном графе с бюджетом, допускающим каждую перестановку, блок оценки выполняет $(n-1)n!$ проверок рёбер. Поэтому трудоёмкость в худшем случае равна $\Theta(n\cdot n!)$. Дополнительная память для текущей перестановки и лучшего маршрута составляет $\Theta(n)$.

В блоке построения маршрута (рисунок~\ref{fig:walk}) на каждом шаге выполняются просмотр $n$ вершин, нормирование списка весов и подготовка дискретного распределения. Последние два действия занимают $O(n)$, поскольку кандидатов не более $n$. За не более чем $n-1$ шагов один муравей выполняет $O(n^2)$ операций. Для завершённых маршрутов просмотр всех вершин выполняется на каждом шаге и даёт $\Theta(n^2)$.

На схеме муравьиного поиска (рисунок~\ref{fig:ants}) в одной итерации построение $K$ маршрутов занимает $O(Kn^2)$, внесение обычных и элитных добавок и сохранение лучшего пути --- $O(Kn+n)$. Инициализация добавок и просмотр матрицы при испарении занимают $\Theta(n^2)$ на итерацию. С учётом начального феромона имеем
\begin{equation}
F_{\text{м}}(n,D,K)=\Theta\!\bigl(n^2+D(n^2+Kn^2)\bigr)
\end{equation}
в худшем случае. Нижняя граница достигается на полном графе при бюджете, достаточном для любого маршрута. При фиксированных $D$ и $K$ это $\Theta(n^2)$. Дополнительная память $\Theta(n^2)$ занята главным образом матрицами $\tau$ и $\Delta$. Муравьи обрабатываются последовательно, поэтому одновременно хранятся только текущий и лучший маршруты.

\subsection*{Вывод}

Перебор проверяет все перестановки и даёт точный результат с факториальной трудоёмкостью. Муравьиный алгоритм выполняет ограниченное число попыток и сохраняет ненулевую вероятность каждого перехода к непосещённому соседу.
